\documentclass[12pt,a4paper]{article}
\usepackage[utf8]{inputenc}
\usepackage[T1]{fontenc}
\usepackage{amsmath,amsfonts,amssymb}
\usepackage{graphicx}
\usepackage{cite}
\usepackage{url}
\usepackage{geometry}
\usepackage{fancyhdr}
\usepackage{setspace}
\usepackage{booktabs}
\usepackage{algorithm}
\usepackage{algorithmic}
\usepackage{xcolor}

% Page setup
\geometry{margin=1in}
\doublespacing
\pagestyle{fancy}
\fancyhf{}
\rhead{\thepage}
\lhead{Online Social-Compliance Verification for AVs}

% Helper: planned (NSFC) content not yet supported by results
\newcommand{\planned}[1]{{\color{blue!60!black}\textit{[PLANNED: #1]}}}
% Helper: figure placeholder box (keeps the file compilable without image files)
\newcommand{\figph}[2]{\fbox{\parbox[c][#1][c]{0.92\linewidth}{\centering\itshape #2}}}

% Title and author information
\title{Online runtime verification of socially compliant autonomous driving}
\author{
    Author Name\thanks{Corresponding author: author@institution.edu} \\
    \textit{Department, Institution} \\
    \textit{City, Country}
}
\date{\today}

\begin{document}

\maketitle

\begin{abstract}
\noindent
Autonomous vehicles are certified for collision safety, yet they operate within human
traffic where competent behaviour is governed by implicit, context-dependent social
norms. We recast social compliance as \emph{state-conditioned online runtime
verification}. From 38{,}228 multi-source human interactions we estimate empirical
normative envelopes of an interaction preference value (IPV) over a
risk\,$\times$\,geometry\,$\times$\,role\,$\times$\,time state space, and build an online
verifier that estimates each agent's IPV causally from past trajectory and scores its
deviation from the state-conditioned envelope in real time. Social compliance is strongly
state-dependent: the priority-related preference gap reverses sign with collision risk
($+0.058$ to $-0.034$) and is robust across data sources and binning choices. The
verifier's social signal is causal and deployable; we explicitly quantify the residual
online/offline gap, which is confined to risk-state recovery rather than to the social
signal. We position the verifier as an auditable monitoring layer that complements, rather
than relabels, collision-safety verification, advancing certifiable, human-aligned autonomy.
\end{abstract}

% ======================================================================
\section{Introduction}
% ----------------------------------------------------------------------
Autonomous vehicles (AVs) must share the road with human drivers, cyclists and
pedestrians whose behaviour is shaped not only by physical constraints but also by
implicit social norms. A manoeuvre that is collision-free in isolation may still be
perceived as aggressive, hesitant or norm-violating when performed in the wrong
interaction context, eroding the predictability on which mixed human--machine traffic
depends. Recent controlled experiments show that machine agents embedded in human groups
can suppress beneficial social norms such as reciprocity, with collective
consequences~\cite{shirado2023reciprocity}. Social competence is therefore not a cosmetic
add-on but part of what it means for autonomy to be safe \emph{in a human world}.

Two established families of methods address adjacent but distinct questions. Offline
behavioural evaluation asks ``did the AV, on average, drive like a human?'', typically by
reporting a single aggregate sociality score over a dataset. Formal safety verification
asks ``will this trajectory collide or violate a hard constraint?'', and can guarantee
legal safety as an online layer around a planner~\cite{pek2020online}. Neither answers the
question a socially aware AV actually faces moment to moment: \emph{given the current
interaction state, is this behaviour within the range that competent humans exhibit?} This
question is local in time, conditional on context, and graded rather than binary.

We formulate socially compliant driving as the \emph{runtime verification of
online-inferred interaction preferences}. The property being monitored is not a
hand-written temporal-logic specification but a \emph{data-driven, state-conditioned
empirical norm} learned from large-scale human interaction. Concretely, we represent
driving norms as a state-conditioned social-norm field: for each interaction state, the
field defines an empirical envelope of socially typical behaviour, expressed through the
interaction preference value (IPV), a trajectory-derived measure of how an agent balances
its own progress against the group's. Given this field, an AV can ask online whether its
current or planned behaviour falls within the empirical envelope expected under the present
state, and a monitor can expose the answer to the planner as a warning, soft cost,
fallback trigger or audit record.

This framing differs from three existing paradigms. Unlike offline AV evaluation, it is
online, per-timestep and state-conditioned rather than aggregate. Unlike runtime
verification of formal specifications, the monitored property is an empirical normative
distribution estimated from human data rather than a logical rule. Unlike socially
compliant planning, it does not encode sociality into a planner's objective but provides a
pluggable, planner-agnostic monitor that any planner can be checked against.

Our contributions are fourfold. First, we formalise social compliance as a state-conditioned
runtime-verification problem and instantiate it as a four-layer online verifier. Second, we
provide strong empirical evidence that social norms are state-dependent: the social meaning
of right-of-way reverses with collision risk, and this reversal is robust across four
independent data sources. Third, we show that the verifier's social signal---a dynamic IPV
estimated causally from past trajectory---is online by construction, and we honestly
quantify the residual runtime gap, which is confined to recovering the risk state online.
Fourth, \planned{external validation on an independent real-vehicle challenge dataset, and
a demonstration that social-compliance deviation carries information about outcomes that
collision-safety checks miss}. The present manuscript reports the first three contributions
in full and scopes the fourth as planned validation.

% ======================================================================
\section{Results}
% ----------------------------------------------------------------------

\subsection{Social compliance is a state, not a score}
We analysed 38{,}228 interaction cases drawn from four naturalistic sources within the
InterHub corpus (Waymo 23{,}218; nuPlan 7{,}499; Lyft 5{,}105; Argoverse-2 2{,}406;
3{,}695{,}981 trajectory frames). For each interacting agent we computed the IPV,
parameterised as an angle $\theta\in[-\pi/2,\pi/2]$, where $\theta>0$ denotes cooperative
or prosocial behaviour and $\theta<0$ competitive behaviour.

The social meaning of right-of-way is not a static label but is gated by collision risk
(Fig.~\ref{fig:state-dependence}). The difference in IPV between the priority and
non-priority agent reverses sign as post-encroachment time (PET) increases: under high risk
(PET\,$\le$\,1.0\,s) the priority agent is markedly more prosocial
($+0.058$, 95\%\,CI $[+0.050,+0.067]$, $n=9{,}799$); at intermediate risk the gap vanishes
($+0.001$, $[-0.004,+0.006]$, $n=22{,}833$); and under low risk (PET\,$>$\,2.0\,s) it
inverts, the priority agent becoming relatively more competitive
($-0.034$, $[-0.045,-0.023]$, $n=4{,}863$). Right-of-way is thus a risk-modulated
negotiation variable rather than a fixed disposition.

The same state dependence appears in the normative envelopes themselves
(Fig.~\ref{fig:envelopes}). Across 16 high-support main-text cells (of 432 total; 79 with
$n\ge200$), the median IPV and prosocial share fall monotonically as risk decreases: from
median $0.137$ and prosocial share $0.738$ in the high-risk priority-conflict cell to
median $-0.003$ and share $0.476$ in the low-risk priority-pre cell. A coarse geometric
prior is stable across sources: merging/proceeding and straight--straight configurations
are prosocial in 4/4 datasets. A hierarchical partial-pooling estimator recovers these
envelopes with leave-one-dataset-out median absolute error $0.142$, below both
dataset-specific ($0.145$) and fully pooled ($0.150$) alternatives.

These conclusions survive extensive robustness checks (Extended Data): the
direction of the risk-gated priority effect is preserved under dropping Waymo, under
three-, five- and quantile-binning of PET, under geometry coarsening, and under
reliability thresholds of $n\ge100/200/300$. Social compliance therefore cannot be reduced
to a single average; it must be described as a function of
risk\,$\times$\,geometry\,$\times$\,role\,$\times$\,time.

\subsection{An online verifier with a causal social signal}
We instantiate the verifier as four layers (Fig.~\ref{fig:architecture}): a state
recogniser, an empirical normative envelope, a deviation scorer and a planner interface.
The design explicitly separates an \emph{offline calibration layer} from an \emph{online
monitoring layer}.

The monitored social signal is online by construction. The dynamic IPV is estimated
causally from a trailing window of observed trajectory for both the ego agent and its
counterpart; it uses no future information and is therefore a strict-online quantity,
distinct from the full-window, case-mean IPV that is used only as an offline calibration
label. What remains genuinely incomplete online is the \emph{risk state} used to index the
envelope. Reconstructing it from causal features, the best online risk proxy---the minimum
of a projected anticipated-PET and a pair-closing time-to-collision---agrees with the
offline PET bin in $52.1\%$ (balanced) to $57.8\%$ (weighted) of cases, with high-risk
recall $0.67$ and precision $0.48$; a rolling phase detector aligns with offline phases on
$72.2\%$ of frames and triggers a median $0.6$\,s before the offline conflict onset, but
overlaps the offline conflict window with an intersection-over-union of only $0.30$. We
report this runtime gap as a concept-level result rather than concealing it: online
information is intrinsically partial, which is precisely why a runtime verifier---rather
than an offline label---is required.

\subsection{State-conditioned envelopes are interpretable monitoring references}
Under a unified leave-one-dataset-out protocol (estimate and calibrate on three sources,
evaluate on the held-out source) we compared the full state-conditioned envelope against
scalar, identity, priority-only, risk-only and geometry-only baselines
(Fig.~\ref{fig:ladder}). The full-state model is consistently but \emph{modestly} better:
relative to the stateless scalar it reduces interval width by $2.1\%$ ($0.684$ vs $0.699$)
and pinball loss by $3.5\%$ ($0.0598$ vs $0.0619$), and raises a synthetic-deviation
detection AUC from $0.797$ to $0.814$ (fold SD\,$\approx0.038$). Crucially, coverage does
not improve ($0.759$ vs $0.768$): the gain reflects sharper, more informative references,
not better-calibrated guarantees. A shuffled-state negative control collapses to scalar
performance ($\approx0.507$), confirming the gain originates in genuine state conditioning.
We therefore position the state-conditioned verifier as an \emph{interpretable monitoring
architecture and normative reference}, and explicitly \emph{not} as a predictively superior
model.

\subsection{A social-deviation signal can precede the conflict window}
Descriptively, a non-zero IPV is already present before the labelled interaction onset in
$71.4\%$ of cases, suggesting that social intent is observable early. However, when
evaluated as a strict online early-warning detector at a controlled $5\%$ false-positive
rate (Fig.~\ref{fig:earlywarning}), the capability is limited: ROC-AUC $0.648$, recall
$12.8\%$ ($[11.9,13.7]$) at lead $\ge0$\,s, $7.4\%$ ($[6.6,8.1]$) at $\ge1$\,s and $1.1\%$
at $\ge2$\,s, with a detected-horizon median of $1.1$\,s ($[0.2,1.9]$); the full-state model
does not exceed the scalar baseline here. We therefore report early warning as a
\emph{proof of concept}: the descriptive ``signal precedes conflict'' observation does not
yet constitute a high-recall online alarm.

\subsection{Verifier output as a planner-facing channel}
The deviation score maps to planner-facing outputs---a one-sided soft cost active only where
the norm expects prosociality, a quality-gated warning, a fallback trigger and a social-risk
monitor record (Fig.~\ref{fig:planner}). Injecting a counterfactual ``more competitive''
violation into $4{,}800$ case-agents raises the fallback-trigger rate from $3.0\%$ to
$12.9\%$ and lifts the median pre-onset cost by $+0.252$, while leaving compliant streams
largely unpenalised (mean social cost $0.134$, false-trigger rate $3.0\%$). This is an
\emph{interface demonstration} of actionability, not a closed-loop planner result: it shows
that verifier output can be exposed as a usable signal, without claiming improved planner
performance or any safety guarantee. Offline deviation sensitivity also transfers
moderately across sources (held-out deviation AUC $0.830$ vs scalar $0.817$), whereas online
calibration remains source-dependent (Waymo held-out coverage $0.706$; Lyft online
false-positive rate $19.8\%$).

\subsection{External validation on an independent real-vehicle challenge}
\planned{We will run the verifier on the award-winning algorithms of a national real-vehicle
autonomous-driving challenge, whose official rankings, per-scenario scores and safety-event
records provide an independent outcome ground truth. The analysis will test (i) criterion
validity---whether algorithms ranked higher sit closer to the human normative envelope; and
(ii) the consequence chain---whether higher social-compliance deviation predicts more
safety events and lower scores (Fig.~\ref{fig:nsfc-validation}).}

\subsection{Discriminant value beyond collision-safety verification}
\planned{Using the same challenge data, we will quantify the verifier's value relative to
conventional safety checks: identifying cases that pass formal safety verification yet are
flagged by the social verifier, and testing whether social-compliance deviation retains
predictive power for independent scores and safety events after controlling for standard
safety and rule-compliance metrics (Fig.~\ref{fig:discriminant}). This establishes the
verifier as complementary to, not a relabelling of, collision-safety verification.}

% ======================================================================
\section{Discussion}
% ----------------------------------------------------------------------
We have reframed AV social behaviour from an offline evaluation metric into an online
runtime-verification property, learned empirical state-conditioned normative envelopes from
multi-source human data, and shown that the social signal monitored by the verifier is
causal and deployable. The central empirical result---that social compliance is strongly
state-dependent, with right-of-way preference reversing sign across the risk axis---argues
that any scalar ``sociality score'' is ill-posed, and that compliance must be judged
relative to a state-conditioned human envelope.

Our claims are deliberately bounded. The verifier is an empirical, probabilistic monitor,
not a formal proof, and it complements rather than replaces safety verification such as
reachability-based legal-safety layers~\cite{pek2020online}. Three boundaries are explicit.
The predictive advantage of state conditioning over a scalar baseline is small and is not
better-calibrated; we therefore frame the contribution as an interpretable monitoring
reference. Early warning is a proof of concept, limited by the absence of causal
conflict-event timestamps and by synthetic violation labels. A further rival explanation is
that the IPV re-encodes kinematics such as relative speed, braking or passing order rather
than a genuine social preference; the shuffled-state control argues against a pure artifact,
but kinematic-only and IPV-ablated comparisons are needed to exclude it. And generalisation is local:
the corpus is Waymo-dominated (61\% of cases, 85\% of AV cases), online calibration drifts
across sources, and only coarse, high-support states transfer reliably. The planned
external validation directly targets these boundaries by supplying real algorithms and
independent outcome labels.

Beyond driving, the framework offers a template for monitoring whether autonomous agents
behave within human normative envelopes online, under uncertainty and distribution shift,
rather than optimising a single average behavioural metric---a question of growing
importance wherever machine agents act within human social systems.

% ======================================================================
\section{Methods}
% ----------------------------------------------------------------------

\subsection{Data and interaction extraction}
We used 38{,}228 interaction cases from four sources within InterHub (Waymo, nuPlan, Lyft,
Argoverse-2), comprising 3{,}695{,}981 trajectory frames; 17 all-zero-IPV cases were
retained as boundary cases. Each case identifies a dyadic interaction with an annotated
priority (right-of-way) agent. We emphasise that the priority label denotes the
right-of-way agent identity, not a numerical rank or the realised passing order.

\subsection{IPV estimation and causal online inference}
The interaction preference value quantifies an agent's balance between self-interest and
group reward, parameterised as $\theta\in[-\pi/2,\pi/2]$. We estimate a \emph{dynamic} IPV
for both agents from a trailing window of observed trajectory, yielding a time-varying,
strictly causal signal suitable for online deployment. \emph{[To specify: window length;
whether predicted-intention trajectories are used in addition to observed history; and the
convergence/bias of the rolling estimate relative to the full-window estimate.]} The
full-window, case-mean IPV is used only as an offline calibration label and is never an
online input.

\subsection{State space and normative envelopes}
The primary state space is
$\text{risk\_bin(PET, 3)}\times\text{geometry(path category)}\times
\text{role(priority/non-priority)}\times\text{phase(pre/conflict/post)}$, giving 432 cells,
of which 79 are reliable ($n\ge200$) and 16 are high-support main-text cells. Envelopes are
estimated as per-state quantiles ($q_{05},q_{50},q_{95}$) via hierarchical partial pooling
with bootstrap intervals; sparse cells fall back to parent levels. To match the online
signal, envelopes should be calibrated on the same dynamic-IPV process at matched
window length / interaction progress, exploiting the existing time axis of the state space,
so that the online comparison is like-for-like rather than rolling-versus-full-window.

\subsection{Online verification: dynamic envelopes and the conditional criterion}

\paragraph{Notation.}
We consider a dyadic interaction between an ego agent $i$ and its principal counterpart
$j$. At time $t$ we observe a trailing window $W_t=\{x(t'):\,t-\Delta\le t'\le t\}$ of
kinematic states. A causal estimator $g$ returns a \emph{dynamic} interaction preference
value (IPV) for each agent, $\hat\theta_i(t)=g(W_t;i)$ and $\hat\theta_j(t)=g(W_t;j)$, with
$\theta\in[-\pi/2,\pi/2]$ and $\theta>0$ denoting prosocial behaviour. Because $g$ reads only
past observations, $\hat\theta(t)$ is a strict-online signal, distinct from the full-window,
case-mean IPV used only as an offline calibration label.
\emph{[TO SPECIFY: window length $\Delta$; whether predicted-intention trajectories augment
$W_t$; reported convergence/bias of $\hat\theta(t)$ relative to the full-window estimate.]}

\paragraph{Interaction state.}
We summarise the situation by a state $s(t)=(r(t),h,\rho_i,\tau(t))$, where $r(t)$ is an
online risk bin from the hybrid proxy $r=\min(\widehat{\mathrm{APET}},\mathrm{TTC}_{\mathrm{close}})$,
$h$ is the coarse geometry (path category), $\rho_i\in\{\text{priority},\text{non-priority}\}$
is the map-derived role, and $\tau(t)\in[0,1]$ is the interaction progress from a causal
rolling-phase detector. The progress index $\tau$ is essential: it makes the norm a function
of how far the interaction has developed, so that the online signal is compared with a norm
estimated at the same stage rather than with a whole-episode summary.

\paragraph{Dynamic normative envelope.}
Offline, for each state cell $c=(r,h,\rho_i)$ and progress $\tau$, we estimate the
distribution of the dynamic IPV computed with the \emph{same} window $\Delta$ on human data,
not the full-window mean. The marginal envelope is the set of conditional quantiles
\begin{equation}
Q_\alpha(c,\tau)=\inf\{\theta:\hat F(\theta\mid c,\tau)\ge\alpha\},\qquad
\text{reported as } [q_{05},q_{50},q_{95}].
\end{equation}
Calibrating on the rolling estimator removes the rolling-versus-full-window mismatch, so an
online value is scored against a like-for-like reference. Sparse cells back off to a parent
level by hierarchical partial pooling.

\paragraph{Conditional criterion $P(\theta_i\mid\theta_j)$.}
Human interactions show event-level coupling: simultaneous strong competition is rare, so an
agent's socially reasonable preference depends on its counterpart's. We therefore verify the
ego against the \emph{conditional} norm. For state $c$ and progress $\tau$ we estimate
$\hat F(\theta_i\mid\theta_j,c,\tau)$ and define the reasonable band
\begin{equation}
B_{1-\eta}(\theta_j,c,\tau)=\bigl[\,Q_{\eta/2}(\theta_j,c,\tau),\;Q_{1-\eta/2}(\theta_j,c,\tau)\,\bigr],
\end{equation}
the central $1-\eta$ conditional interval of ego IPV given the counterpart's current IPV.
When the counterpart is competitive ($\theta_j\ll0$) the band shifts prosocial; an ego that
is simultaneously competitive then falls outside the band, flagging the configuration that
empirically precedes conflict.
\emph{[TO SPECIFY: conditional quantiles via binning on $\theta_j$ within each cell, or a
monotone quantile regression; chosen per cell support.]}

\paragraph{Deviation score.}
The signed, normalised nonconformity of the ego is
\begin{equation}
\delta_i(t)=\frac{\hat\theta_i(t)-Q_{1/2}\!\bigl(\hat\theta_j(t),c,\tau\bigr)}
{w\!\bigl(\hat\theta_j(t),c,\tau\bigr)},
\qquad
w=\tfrac{1}{2}\bigl(Q_{1-\eta/2}-Q_{\eta/2}\bigr),
\end{equation}
so $|\delta_i(t)|>1$ marks a value outside the reasonable band and the sign gives the
direction (competitive versus over-yielding). For a calibrated flag we wrap $\delta$ in a
per-cell (Mondrian) split-conformal threshold $\kappa_c$ estimated on held-out human data; a
flag at level $\beta$ then means $\delta$ exceeds the $1-\beta$ human quantile for that cell.
Because coverage degrades under source shift, $\kappa_c$ is calibrated source-aware and
reported with its empirical coverage rather than assumed nominal. The online score is
smoothed by an exponentially weighted moving average,
$\bar\delta(t)=\lambda\,\delta_i(t)+(1-\lambda)\,\bar\delta(t-1)$, and persistence is required
before escalation.

\paragraph{Verdict and planner interface.}
A verdict is emitted only after quality gates pass: sufficient cell support, admissible input
provenance (leakage contract), bounded estimator uncertainty, and temporal hysteresis. The
gated score maps to a graded output: a monitor record always; a one-sided soft cost
$c_{\mathrm{social}}=\max(0,\bar\delta-\theta_{\mathrm{warn}})\cdot\mathbf{1}[\text{norm prosocial}]$;
a warning above $\theta_{\mathrm{warn}}$; and a fallback candidate above $\theta_{\mathrm{fb}}$
under sustained high risk.
\emph{[TO SPECIFY: thresholds $\theta_{\mathrm{warn}},\theta_{\mathrm{fb}}$ set on the
calibration set; demonstrative values used in pilot analyses were $0.25$ and $0.85$.]}

\begin{algorithm}[t]
\caption{Offline calibration of dynamic and conditional envelopes}
\label{alg:calib}
\begin{algorithmic}[1]
\REQUIRE human interactions $\mathcal{D}$; window $\Delta$; levels $\alpha,\eta,\beta$;
         min support $n_{\min}$
\ENSURE marginal $Q$, conditional bands $B$, per-cell thresholds $\kappa$
\STATE split $\mathcal{D}$ into estimation set $\mathcal{D}_e$ and calibration set $\mathcal{D}_k$
\STATE $\mathcal{T}\leftarrow\emptyset$
\FOR{each interaction in $\mathcal{D}_e$ and each time $t$}
    \STATE $\theta_i,\theta_j\leftarrow g(W_t;i),\,g(W_t;j)$ \COMMENT{causal dynamic IPV}
    \STATE $(r,h,\rho_i,\tau)\leftarrow \textsc{RecognizeState}(W_t)$
    \STATE append $\bigl((r,h,\rho_i),\tau,\theta_i,\theta_j\bigr)$ to $\mathcal{T}$
\ENDFOR
\FOR{each cell $c$ and progress bin $\tau$ in $\mathcal{T}$}
    \IF{$\mathrm{support}(c,\tau)<n_{\min}$}
        \STATE attach $c$ to $\textsc{Parent}(c)$ \COMMENT{hierarchical fallback}
    \ENDIF
    \STATE $Q_\alpha(c,\tau)\leftarrow$ empirical quantiles of $\theta_i$
    \STATE $B_{1-\eta}(\theta_j,c,\tau)\leftarrow$ conditional quantiles of $\theta_i$ given $\theta_j$
\ENDFOR
\STATE on $\mathcal{D}_k$: $\kappa_c\leftarrow (1-\beta)$ quantile of $|\delta|$ per cell \COMMENT{Mondrian split-conformal, source-aware}
\RETURN $Q,\,B,\,\kappa$
\end{algorithmic}
\end{algorithm}

\begin{algorithm}[t]
\caption{Online social-compliance verification step at time $t$}
\label{alg:online}
\begin{algorithmic}[1]
\REQUIRE window $W_t$; calibrated $Q,B,\kappa$; gates; previous EWMA $\bar\delta_{t-1}$;
         smoothing $\lambda$; thresholds $\theta_{\mathrm{warn}},\theta_{\mathrm{fb}}$
\ENSURE verdict and planner-facing signal
\STATE $\theta_i,\theta_j\leftarrow g(W_t;i),\,g(W_t;j)$ \COMMENT{strict-online; no future info}
\STATE $(r,h,\rho_i,\tau)\leftarrow \textsc{RecognizeState}(W_t)$
\STATE $c\leftarrow(r,h,\rho_i)$;\quad \textbf{if} $\mathrm{support}(c,\tau)<n_{\min}$ \textbf{then} $c\leftarrow\textsc{Parent}(c)$
\STATE $m\leftarrow Q_{1/2}(\theta_j,c,\tau)$;\quad $w\leftarrow\tfrac12\bigl(Q_{1-\eta/2}-Q_{\eta/2}\bigr)(\theta_j,c,\tau)$
\STATE $\delta\leftarrow (\theta_i-m)/w$ \COMMENT{signed conditional nonconformity}
\STATE $\bar\delta\leftarrow \lambda\,\delta+(1-\lambda)\,\bar\delta_{t-1}$ \COMMENT{EWMA}
\IF{\textbf{not} \textsc{QualityGates}$(c,\text{provenance},\text{uncertainty})$}
    \RETURN $\textsc{Monitor}(\bar\delta)$ \COMMENT{insufficient support/quality: record only}
\ENDIF
\IF{$\bar\delta>\kappa_c$ \textbf{and} $r$ high \textbf{and} persistent}
    \STATE $\text{signal}\leftarrow\textsc{FallbackCandidate}$
\ELSIF{$\bar\delta>\theta_{\mathrm{warn}}$}
    \STATE $\text{signal}\leftarrow\textsc{Warning}$, with soft cost $\max(0,\bar\delta-\theta_{\mathrm{warn}})\cdot\mathbf{1}[\text{norm prosocial}]$
\ELSE
    \STATE $\text{signal}\leftarrow\textsc{Monitor}$
\ENDIF
\RETURN $(\bar\delta,\ \text{signal})$
\end{algorithmic}
\end{algorithm}

\noindent\textbf{Complexity and determinism.} Each online step is $O(1)$ given precomputed
envelope tables (state recognition, two table look-ups and a scalar update), so the verifier
runs within a planning cycle and is deterministic given $W_t$, supporting auditability.

\subsection{Leakage contract}
We enforce a strict separation of online and offline variables. Dynamic/rolling IPV and
current kinematics are strict-online and admissible as verifier inputs; anticipated-PET/TTC
proxies and map-derived priority are conditional-online; observed PET, realised passing
order, full-window IPV, closest-frame and post-hoc phase labels are offline-only and are
used solely for calibration and evaluation, never as deployed inputs.

\subsection{Evaluation protocols}
Model comparison uses leave-one-dataset-out estimation and evaluation across the four
sources, scoring coverage, interval width, pinball loss, median absolute error and a
synthetic-deviation detection AUC, with shuffled-state and random cross-state negative
controls. Planned ablations include kinematic-only and IPV-removed variants that test
whether the social signal is separable from raw kinematics. Early warning is evaluated at a
controlled $5\%$ false-positive rate with
Wilson confidence intervals and lead-time-binned recall. The planner interface is evaluated
by counterfactual violation injection with an over-conservatism check on compliant streams.

\subsection{Planned external validation protocol}
\planned{The verifier will be applied to the trajectories of award-winning algorithms from a
national real-vehicle challenge. We will align scenarios, compute per-algorithm and
per-scenario deviation, and relate deviation to official rankings, per-scenario scores and
safety events via correlation and incremental regression that controls for conventional
safety and rule-compliance metrics; where human reference trajectories are available, AV
algorithms and humans will be placed within the same normative envelope.}

% ======================================================================
% Figures (placeholders kept compilable; replace \figph with \includegraphics)
% ======================================================================
\begin{figure}[t]\centering
\figph{3.2cm}{Figure 1 placeholder. Source: agent\_F\_figures/round5/fig1\_architecture}
\caption{\textbf{Online social-compliance verification framework.} An offline calibration
layer maps multi-source interactions to a state space and empirical normative envelopes
(PET, start/end and full-window IPV used only as offline labels), and an online monitoring
layer maps streaming features (causal risk proxy, rolling phase, geometry, role, dynamic
IPV) through a state recogniser, envelope lookup with fallback, deviation scorer and planner
interface. The residual runtime gap is annotated, not hidden.}
\label{fig:architecture}
\end{figure}

\begin{figure}[t]\centering
\figph{3.2cm}{Figure 2 placeholder. Source: agent\_F\_figures/round2/fig2\_state\_space\_map}
\caption{\textbf{Empirical social-compliance state space.} Support and sample size of each
risk\,$\times$\,geometry\,$\times$\,role\,$\times$\,phase cell, highlighting the 16
high-support main-text cells and the sparse/Extended-Data regions.}
\label{fig:state-space}
\end{figure}

\begin{figure}[t]\centering
\figph{3.2cm}{Figure 3 placeholder. Source: agent\_F\_figures/round1\&2 (PET-gating, envelopes, geometry)}
\caption{\textbf{Social norms are state-dependent.} The priority-minus-non-priority IPV gap
reverses with PET risk ($+0.058\!\to\!+0.001\!\to\!-0.034$) with 95\% confidence intervals
and per-source forest plots; coarse geometry provides a cross-source prosocial prior;
state-conditioned normative envelopes (median, 25--75\%, 10--90\%) flatten as risk decreases.}
\label{fig:envelopes}
\end{figure}

\begin{figure}[t]\centering
\figph{3.0cm}{Figure 4 placeholder. Source: agent\_F\_figures/round2\&4 (online feasibility, deviation over time)}
\caption{\textbf{Causal online verification and the runtime gap.} The dynamic IPV deviation
is computed online from past trajectory; the online risk proxy recovers $\sim$55\% of the
offline PET stratification and the rolling phase detector overlaps the offline conflict
window at IoU $0.30$. The gap is confined to risk-state indexing, not to the social signal.}
\label{fig:state-dependence}
\end{figure}

\begin{figure}[t]\centering
\figph{3.0cm}{Figure 5 placeholder. Source: agent\_F\_figures/round3/fig6\_baseline\_ladder}
\caption{\textbf{State-conditioned envelope versus baselines (leave-one-dataset-out).}
Coverage and width, pinball and median absolute error, and synthetic-deviation AUC. The
full-state model is sharper and has the highest AUC ($0.814$ vs $0.797$) but does not
improve coverage; gains are modest and originate in state conditioning (shuffled-state
control $\approx$ scalar).}
\label{fig:ladder}
\end{figure}

\begin{figure}[t]\centering
\figph{3.0cm}{Figure 6 placeholder. Source: agent\_F\_figures/round4/fig5\_early\_warning}
\caption{\textbf{Early-warning performance (proof of concept).} At a controlled $5\%$
false-positive rate, ROC-AUC is $0.648$ and recall is $7.4\%$ at lead $\ge1$\,s; the
full-state model does not exceed the scalar baseline. Suitable as a planner soft-cost or
fallback signal, not as a reliable low-false-alarm alarm.}
\label{fig:earlywarning}
\end{figure}

\begin{figure}[t]\centering
\figph{3.0cm}{Figure 7 placeholder. Source: agent\_F\_figures/round5/fig7\_planner\_demo}
\caption{\textbf{Planner-facing soft-cost demonstration.} Counterfactual violation injection
raises the fallback-trigger rate from $3.0\%$ to $12.9\%$ and the median pre-onset cost by
$+0.252$, while leaving compliant streams largely unpenalised. An interface demonstration,
not a closed-loop or safety guarantee.}
\label{fig:planner}
\end{figure}

\begin{figure}[t]\centering
\figph{2.8cm}{Figure 8 placeholder \planned{NSFC external validation}}
\caption{\planned{\textbf{External validation on an independent real-vehicle challenge.}
Per-algorithm social-compliance deviation versus official ranking and per-scenario scores;
deviation versus safety-event rate.}}
\label{fig:nsfc-validation}
\end{figure}

\begin{figure}[t]\centering
\figph{2.8cm}{Figure 9 placeholder \planned{NSFC discriminant value}}
\caption{\planned{\textbf{Discriminant value beyond safety verification.} Cases that pass
safety verification yet are flagged by the social verifier, and incremental prediction of
scores/events after controlling for conventional safety metrics.}}
\label{fig:discriminant}
\end{figure}

% ======================================================================
\section*{Data Availability}
InterHub-derived interaction cases and source datasets (Waymo, nuPlan, Lyft, Argoverse-2)
are governed by their respective licences; processed state-space derivatives will be made
available upon publication. \planned{Real-vehicle challenge data availability to be stated
per the challenge organiser's terms.}

\section*{Code Availability}
Implementation, leakage-contract unit tests and figure-generation code will be released
upon publication.

\section*{References}
\bibliographystyle{unsrt}
\bibliography{bibliography/biblio}

\section*{Acknowledgements}
% (Funding sources, institutional support, data providers)

\section*{Author Contributions}
% (Individual contributions following Nature journal format)

\section*{Competing Interests}
The authors declare no competing interests.

\section*{Additional Information}
\subsection*{Supplementary Information}
\subsection*{Supplementary Tables}
\subsection*{Supplementary Figures}

\end{document}
